\documentclass{article}

\usepackage[a4paper]{geometry}
\usepackage[english]{babel}
\usepackage{parskip}

\usepackage{amsmath, amssymb}
\usepackage{graphicx}
\usepackage{subcaption}
\usepackage{float}
\usepackage[hidelinks]{hyperref}

\usepackage{listings}
\usepackage{xcolor}

\definecolor{codegreen}{rgb}{0,0.6,0}
\definecolor{codegray}{rgb}{0.5,0.5,0.5}
\definecolor{codepurple}{rgb}{0.58,0,0.82}
\definecolor{backcolour}{rgb}{0.95,0.95,0.92}

\lstdefinestyle{mystyle}{
  backgroundcolor=\color{backcolour},
  commentstyle=\color{codegreen},
  keywordstyle=\color{magenta},
  numberstyle=\tiny\color{codegray},
  stringstyle=\color{codepurple},
  basicstyle=\ttfamily\footnotesize,
  breaklines=true,
  captionpos=b,
  keepspaces=true,
  numbers=left,
  numbersep=5pt,
  showstringspaces=false,
  tabsize=2
}
\lstset{style=mystyle}

\newcommand{\CodeDir}{Code}
\newcommand{\PlotDir}{Plots}
\newcommand{\CalcDir}{Calculations}

\title{NURa Hand-in 2}
\author{Niels de Klerk (s3640477)}
\date{March 19, 2026}

\begin{document}
\maketitle


\section{Satellite galaxies around a massive central}

\subsection{(1a) Determining the normalization}
The satellite galaxies around a central follow the number density
\begin{equation}
    n(x) = A\langle N_{sat}\rangle\left(\frac{x}{b}\right)^{a-3}\exp\left[-\left(\frac{x}{b}\right)^c\right]
\end{equation}

where $x$ is the relative radius to the virial radius $(x\equiv r/r_\mathrm{vir})$. This distribution has to normalize to $\langle N_{sat}\rangle$ when integrating from $x_\mathrm{min}=0$ to $x_\mathrm{max}=5$. 

\begin{equation}
    \int dV n(x) = \langle N_{sat}\rangle
\end{equation}

Rewriting $dV$ in spherical coordinates, we find that $dV = 4\pi x^2$. Therefore, the integral that determines the value for $A$ is

\begin{equation}
    \int_{x=0}^{x=5} dx\; 4\pi x^2n(x) = \langle N_{sat}\rangle
\end{equation}

Filling in the equation for $n(x)$ and isolating $A$ results in
\begin{equation}\label{eq:normalization constant}
    A = \frac{\langle N_{sat}\rangle}{\int_{x=0}^{x=5} dx\; 4\pi x^2 \langle N_{sat}\rangle\left(\frac{x}{b}\right)^{a-3}\exp\left[-\left(\frac{x}{b}\right)^c\right]}
\end{equation}

To evaluate this integral, we will use Romberg integration. This numerical integration technique combines different estimates of the integral $\int f(x)dx$ into a better approximation. It starts by finding the area between the two endpoints of the integral $r_0$. For the next estimate, it uses the previous estimate and the sum of function evaluations between the endpoints, scaled by the step size $r_{i+1}=\frac{1}{2}\left[r_i+\Delta_{i+1}\sum_j f (x_j)\right]$. 

After repeating this step $m-1$ times, there are $m$ estimates for the value of the integral. By using the equation 
\begin{equation}
    r_{i,j+1}=\frac{4^{j+1}{r_{i+1, j}-r_{i,j}}}{4^{j+1}-1}
\end{equation}
where the initial values are $r_{i,0}$. This step repeats until a single value remains, obtained by combining all the starting estimates. 

Applying the Romberg integration to \autoref{eq:normalization constant}, we find the following value for the normalization constant.
\begin{table}[H]
\centering
\caption{Normalization constant of the satellite galaxy profile.}
\begin{tabular}{|c|c|c|}
\hline
Normalization constant \(A\) & Absolute error & Absolute error $\langle N_{sat}\rangle$ \\
\hline
\input{\CalcDir/satellite_A.txt} \\
\hline
\end{tabular}
\end{table}

\subsection{Code for the integrator}

\lstinputlisting[language=Python, firstline=0, lastline=98]{\CodeDir/satellites_integrator_code.txt}

\subsection{(1b) Sampling the satellite distribution}
We would like to sample random galaxies from the number density. To find the probability density, we need to find the number of galaxies between $x$ and $x+dx$. 
\begin{equation}
    \int dx\;N(x) = \int dx\;4\pi x^2n(x)
\end{equation}
So we get $N(x)dx=4\pi x^2n(x)dx$. The probability density needs to be normalized so we get
\begin{equation}
    p(x)dx = \frac{N(x)dx}{\langle N_{sat}\rangle} = \frac{4\pi x^2n(x)dx}{\langle N_{sat}\rangle}
\end{equation}
To sample from this distribution, we will use rejection sampling. For this, we need a random number generator (RNG). The RNG we have implemented has 2 subgenerators: Multiply with Carry (MWC) with base $2^{32}$ and a 64-bit XOR-shifter. The MWC starts with a 64-bit unsigned integer x between 0 and $2^{32}$. Then it applies $a * (x \& (2^{32} - 1)) + (x >> 32)$ and uses only the lowest 32 bits for the random number that is fed into the 64-bit XOR-shifter, which shift the bits left, right, and left by certain amounts to generate a new random number, which is returned. To test the RNG, we calculate the mean, variance, and KS statistic and p-value from the KS test, which should respectively be $\sim0.5$, $\sim1/12$, $\sim0$, and $\gtrsim0.05$. When we generate 1.000.000 samples using the RNG, we do indeed find values that closely follow the expected values, which indicates that this is a good RNG. Now that we have an RNG, we can apply the rejection sampling to the distribution. First, the probability density should be divided by its maximum such that the new maximum is 1. Next, a random x value $x_\mathrm{rand}$ is sampled from a uniform distribution from 0 to 5. After that, a random y value $y_\mathrm{rand}$ is sampled from a uniform distribution from 0 to 1. If $p(x_\mathrm{rand}) > y_\mathrm{rand}$, the $x_\mathrm{rand}$ is accepted into the sample. Otherwise, new x and y values are sampled until the required number of samples is reached.

Using this technique, 10000 samples are drawn from the probability density. Next, the samples are binned in logarithmic bins and scaled by the bin width and the number of samples.

\begin{figure}[H]
\centering
\includegraphics[width=0.85\textwidth]{\PlotDir/my_solution_1b.png}
\caption{Comparison between the analytical satellite distribution \(N(x)\) and the histogram of 10 000 sampled satellite galaxies in logarithmic bins.}
\end{figure}

The sampled distribution largely agrees with the analytical prediction. However, below $\sim 10^{-2}$ there are no samples, whereas the analytical prediction is non-zero. This is due to the range of the relative radius values. There is a small probability that a uniformly sampled x value falls inside this range. Additionally, the probability density in this range is also small, such that even when a uniformly sampled x value falls inside this range, the change that is accepted into the sample is also small. 

\subsection{Code for sampling}

\lstinputlisting[language=Python, firstline=101, lastline=191]{\CodeDir/satellites_sampling_code.txt}
 
\subsection{(1c) Selecting and sorting satellites}
We want to select 100 galaxies from our sample of 10000 galaxies with equal probability, without duplicates, and without rejection. To do this, we randomly shuffle the galaxies and take the first 100. This ensures the conditions. First, an array of 10000 random floats and an index array of elements 0, 1, 2, ..., 10000-1 are created. Next, the array with random numbers is sorted by swapping different elements. When the array with random numbers is sorted, we obtain a randomized index array. The first 100 random indices are used to select galaxies. The relative radii are then sorted to find the number of galaxies within different relative radii. To sort both the array with random numbers and the selection of galaxies, merge sort is used. This algorithm splits up the array into single elements and then combines adjacent elements whilst sorting them until a fully sorted array is recovered. 

\begin{figure}[H]
\centering
\includegraphics[width=0.8\textwidth]{\PlotDir/my_solution_1c.png}
\caption{Cumulative number of selected galaxies as a function of relative radius.}
\end{figure}

This figure shows that most of the galaxies are in a small range of relative radii between 0.02 and 1. 

\subsection{Code for selection and sorting}

\lstinputlisting[language=Python, firstline=194, lastline=292]{\CodeDir/satellites_selection_code.txt}

\subsection{(1d) Numerical derivative}
The analytical formula for the derivative of the number density is given by

\begin{equation}
 \frac{dn}{dx}=\frac{A \langle N_{sat}\rangle}{b}
  \left(\frac{x}{b}\right)^{a - 4}
  \left( c \left(\frac{x}{b}\right)^c - a + 3 \right)
  \exp\left(-\left(\frac{x}{b}\right)^c\right).
\end{equation}
We want to compare the value of the derivative obtained from the analytical formula with that obtained from the Ridders algorithm. Ridders algorithm starts with a list of $m$ estimates for different values of $h$ using the central difference 

\begin{equation}
    \frac{df}{dx}\approx\frac{f(x+h)-f(x-h)}{2h}
\end{equation}

Each consecutive $h$ is reduced by a factor $d$. Similar to Romberg integration, the estimates are combined using
\begin{equation}
    D_{i,j+1}=\frac{d^{2(j+1)}{D_{i+1, j}-D_{i,j}}}{d^{2(j+1)}-1}
\end{equation}
where the initial values are $D_{i,0}$. This step is repeated until one of the following conditions is met:
\begin{enumerate}
    \item The improvement over the previous value is smaller than the target error
    \item The error grows
    \item $m$ iterations are completed
\end{enumerate}
Using this algorithm for the derivative at $x=1$ gives

\begin{table}[H]
\centering
\caption{Derivative of the satellite number density at \(x=1\).}
\begin{tabular}{|c|c|}
\hline
Method & Value \\
\hline
Analytical & \input{\CalcDir/satellite_deriv_analytic.txt} \\
Numerical & \input{\CalcDir/satellite_deriv_numeric.txt} \\
\hline
\end{tabular}
\end{table}

The analytical and numerical derivatives are the same. This is because Ridders' algorithm is an efficient algorithm that can get very low relative errors.

\subsection{Code for derivative calculation}

\lstinputlisting[language=Python, firstline=298]{\CodeDir/satellites_derivative_code.txt}

\section{Heating and cooling in HII regions}


\subsection{(2a) Photoionization heating and radiative recombination cooling}

\textbf{Describe the implementation of your root-finding method (e.g. bisection, Newton-Raphson).}
We want to find the equilibrium temperature of an HII region in Orion, where the only heating is produced by photoionization 
\begin{equation}
    \Gamma_\mathrm{pe} = \alpha_\mathrm{B}n_\mathrm{H}n_\mathrm{e}\psi k_\mathrm{B}T_\mathrm{c}
\end{equation}
and the only cooling is caused by radiative recombination.
\begin{equation}
    \Lambda_\mathrm{rr} = \alpha_\mathrm{B}n_\mathrm{H}n_\mathrm{e}\langle E_\mathrm{rr}\rangle
\end{equation}
with
\begin{equation}
    \langle E_\mathrm{rr}\rangle = \left[0.684 - 0.0416\ln\left(\frac{T_4}{Z^2}\right)\right]k_\mathrm{B}T
\end{equation}

The equilibrium temperature is the temperature where the heating and cooling terms are equal (i.e. $\Gamma_\mathrm{pe}-\Lambda_\mathrm{rr}=0$). To find this temperature, we use a combination of two algorithms. First, we use the false-position method.  It starts with a bracket $[a, b]$. From these two values, a new point $c$ is calculated using 
\begin{equation}
    c = b - \frac{b-a}{f(b)-f(a)}f(b)
\end{equation}
where $f$ is the difference between the cooling and heating equations $\Gamma_\mathrm{pe}-\Lambda_\mathrm{rr}$. Using $c$ and either $a$ or $b$, a new bracket $[c, d]$ is formed, where we demand that $f(c)f(d)<0$ with $d=a$ or $d=b$. This is repeated for some prespecified number of steps before the algorithm switches to the secant method. This method is very similar to the false-position method. The only difference is that the new bracket is always $[b, c]$. False-position converges more slowly than the secant method, but it does not diverge, unlike the secant method. By combining both algorithms this way, we first converge to a better bracket without diverging and then use the quick convergence of the secant method to find the root.
In this case, we switch algorithms after 1 iteration. 
\begin{table}[H]
\centering
\caption{Equilibrium temperature for the simplified heating–cooling model.}
\begin{tabular}{|c|c|c|c|c|}
\hline
Temperature [K] & Absolute error [K] & Relative error & Steps & Time [s]\\
\hline
\input{\CalcDir/equilibrium_temp_simple.txt}
\\
\hline
\end{tabular}
\end{table}

It took 8 iterations to reach an accurate result. The combination of the two algorithms does make it converge much faster, as when only false-position is used, the errors are around a factor $1000$ worse after double the number of steps. 

\subsection{(2b) More realistic model}
To make a more realistic model, we add one extra cooling mechanism, the free-free emission given by
\begin{equation}
    \Lambda_\mathrm{FF} = 0.54T_4^{0.37}\alpha_\mathrm{B}n_\mathrm{H}n_\mathrm{e}k_\mathrm{B}T
\end{equation}
and two heating mechanisms: heating by cosmic rays
\begin{equation}
    \Gamma_\mathrm{CR}=An_\mathrm{e}\xi_\mathrm{CR}
\end{equation}
and heating by MHD
\begin{equation}
    \Gamma_\mathrm{MHD}=8.9\cdot10^{-26}n_\mathrm{H}T_4
\end{equation}
So the new equation we want to solve for the equilibrium temperature is $\Gamma_\mathrm{pe} + \Gamma_\mathrm{CR} + \Gamma_\mathrm{MHD} -\Lambda_\mathrm{rr} - \Lambda_\mathrm{FF}=0$. We again use the same algorithm, but this time switch method after 31 steps. 


\begin{table}[H]
\centering
\caption{Equilibrium temperature for different gas densities.}
\begin{tabular}{|c|c|c|c|c|c|}
\hline
Density \(n_H\) [cm\(^{-3}\)] & $T_\mathrm{eq}$ [K] & Abs. error [K] & Rel. error & Steps & Time [s] \\
\hline
$10^{-4}$ & \input{\CalcDir/equilibrium_low_density.txt} \\
$1$ & \input{\CalcDir/equilibrium_mid_density.txt} \\
$10^4$ & \input{\CalcDir/equilibrium_high_density.txt} \\
\hline
\end{tabular}
\end{table}

The equilibrium temperature of the low gas density is orders of magnitude higher than the equilibrium temperature of the intermediate and high gas densities. Furthermore, it took many more steps to converge to the equilibrium temperature in the more realistic model compared to the simpler one.


\subsection{Code for Exercise 2}

\lstinputlisting[language=Python]{\CodeDir/heatingcooling_code.txt} 

\section{Conclusions}

\end{document}
